\chapter{Writing it up}\label{ch:writeup}

\noindent The last skill is the one nobody teaches: reporting the result so that another
human---including future you---can check it. The memo is not the marketing of the
research; it is the deliverable. This chapter is the contract: claim, evidence, method,
weaknesses, risks, reproduction.

% from-zero:begin
\begin{rulebox}[Before you start]
Chapters~8--14. \emph{Recap:} research that cannot be read, checked and rerun by someone else is not finished.
\end{rulebox}
\begin{wins}[After this chapter you can]
\begin{itemize}
\item structure a memo;
\item back every claim with evidence;
\item make your work reproducible;
\item red-team your own result.
\end{itemize}
\end{wins}
\begin{principle}[How to think about this]
Research is only useful if someone else can read it, check it and rerun it. Write the claim first, show the evidence next, and state what would change your mind. Try to break your own result before anyone else does.
\end{principle}
% from-zero:end

\section{The memo skeleton}\label{sec:memo}

Six sections cover every honest report: what is claimed, the evidence, what was tried
(including the silent trials), what could kill it, the risks, and how to reproduce.
Everything else is style.

\begin{figure}[htbp]\centering
\begin{tikzpicture}[font=\scriptsize]
  \foreach \y/\s/\d in {4.0/1. Claim/{one line: what is true},3.3/2. Evidence/{data + code + outputs},2.6/3. Method/{what was tried, incl.\ $N$},1.9/4. Weaknesses/{what could kill it},1.2/5. Risks/{venue/regime/capacity},0.5/6. Repro/{seeds, pins, how to rerun}}{
    \fill[accent!20,draw=accent] (0,\y) rectangle (7.2,\y+0.55);
    \node[right] at (0.15,\y+0.27) {\textbf{\s}};
    \node[right,text=slate] at (3.4,\y+0.27) {\d};}
\end{tikzpicture}

\caption{The memo skeleton: six sections, each answering a question a reviewer will ask.}
\label{fig:memo}
\end{figure}

\begin{example}\label{ex:memo}
\textbf{Why the skeleton is ordered this way.} Claim first because nobody reads further;
evidence second because the claim is worthless without it; trials third because the
evidence is meaningless without the count.
\begin{solution}
The order mirrors a hostile reviewer's path: What? Prove it. How many tries? Why might
it die? What breaks it? Can I rerun it? A memo that answers in that order survives the
review it was written for.
\end{solution}
\end{example}

\section{Claims carry evidence}\label{sec:claim}

\texttt{report\_check.py} lints a memo for the nine fields this book uses:
\begin{lstlisting}[numbers=none]
memo check: 9/9 fields present
missing: none
the checklist is the deliverable -- a memo without fields is a story
\end{lstlisting}

\begin{example}\label{ex:claim}
\textbf{The claim-evidence pair.} ``Funding carry exists'' is a claim;
``$+0.113$bp/h over 500h, 85\% at the interest floor'' is the evidence.
\begin{solution}
Every sentence that asserts something needs a number attached that another person could
recompute. ``SMA works'' paired with ``$+7.9\%$ vs B\&H $+22.7\%$'' is honest even when
the answer is ``it underperformed''---the evidence column is allowed to say the claim
is wrong.
\end{solution}
\end{example}

\begin{figure}[htbp]\centering
\begin{tikzpicture}[font=\scriptsize]
  \foreach \x/\c/\e in {0.3/{``funding carry exists''}/{+0.113bp/h, 500h, floor 85\%}, 4.0/{``SMA works''}/{+7.9\% 721d, vs B\&H +22.7\%}, 7.7/{``edge in microstructure''}/{IS +1.04, OOS $-0.19$}}{
    \fill[slate!15,draw=slate] (\x,0) rectangle (\x+3.0,0.8) node[midway]{\c};
    \fill[accent!25,draw=accent] (\x,-0.9) rectangle (\x+3.0,-0.1) node[midway]{\e};
    \draw[->,ink] (\x+1.5,0) -- (\x+1.5,-0.1);}
  \node[below,text=slate] at (5.4,-1.0) {every claim carries its evidence or it is marketing};
\end{tikzpicture}

\caption{Every claim maps to evidence---including claims the evidence refutes.}
\label{fig:claim}
\end{figure}

\section{Reproducibility}\label{sec:repro}

\texttt{repro\_demo.py} shows the seed contract:
\begin{lstlisting}[numbers=none]
seed 7: ann ret -31.1%  vol 29.6%
seed 7: ann ret -31.1%  vol 29.6%
seed 99: ann ret +24.5%  vol 32.2%
same seed -> same numbers; pin the seed, pin the data, pin the code
\end{lstlisting}

\begin{example}\label{ex:repro}
\textbf{The three pins.} Data (hash), code (commit), seeds (recorded)---any two without
the third is decorative.
\begin{solution}
Reproducibility is the difference between ``I observed'' and ``I estimate''. A result
that depends on an unpinned random seed is not a result; it is one sample of a
distribution you have not described. The pin discipline from ch07 is what turns the
memo from a story into a measurement.
\end{solution}
\end{example}

\begin{figure}[htbp]\centering
\begin{tikzpicture}[font=\scriptsize,every node/.style={align=center}]
  \node[draw=slate,fill=white,rounded corners=1pt,minimum width=2.3cm] (d) at (0,1.6) {pinned data\\ (hash)};
  \node[draw=slate,fill=white,rounded corners=1pt,minimum width=2.3cm] (c) at (3.0,1.6) {pinned code\\ (commit)};
  \node[draw=slate,fill=white,rounded corners=1pt,minimum width=2.3cm] (s) at (6.0,1.6) {pinned seeds\\ (recorded)};
  \node[draw=ink,fill=goodbg,rounded corners=1pt,minimum width=2.3cm] (o) at (3.0,0.3) {same output\\ $\Rightarrow$ same claim};
  \draw[->,ink] (d.south)--(o.north west);
  \draw[->,ink] (c.south)--(o.north);
  \draw[->,ink] (s.south)--(o.north east);
\end{tikzpicture}

\caption{The reproducibility contract: three pins, one reproducible claim.}
\label{fig:repro}
\end{figure}

\section{Red-teaming your own work}\label{sec:redteam}

\texttt{redteam\_demo.py} runs the attack checklist on this book's SMA result:
\begin{lstlisting}[numbers=none]
  lookahead?     signal uses close_i, fill open_i+1  PASS
  costs?         5bp fee + 2bp slip modeled          PASS
  n trials?      1 window pair, no search            PASS
  regime?        OOS fold -0.52 in quiet regime      FLAG: gate on vol
  capacity?      BTC liquid; footprint <<1% ADV      PASS
4 pass / 1 flag -- the flag is the research, not the failure
\end{lstlisting}

\begin{example}\label{ex:redteam}
\textbf{The flag is the finding.} Four passes, one flag: the quiet-regime OOS fold
loses $-0.52$.
\begin{solution}
The flag is not a reason to kill the memo---it is the memo's most honest section. A
result that reports its own regime boundary is more valuable than one that hides it:
the next researcher knows where to look, and the risk model knows when the signal
is off.
\end{solution}
\end{example}

\begin{figure}[htbp]\centering
\begin{tikzpicture}[font=\scriptsize,every node/.style={draw=slate,fill=white,rounded corners=1pt,align=center}]
  \node (a) at (0,2.0) {read the\\ claim};
  \node (b) [right=1.2cm of a] {check the\\ evidence};
  \node (c) [right=1.2cm of b] {count the\\ trials};
  \node (d) [right=1.2cm of c] {find the\\ flag};
  \draw[->,ink] (a)--(b); \draw[->,ink] (b)--(c); \draw[->,ink] (c)--(d);
  \node[below=0.4cm of b,draw=none,text=slate] {a reviewer's path is the memo's outline};
\end{tikzpicture}

\caption{The reviewer's path: claim, evidence, trial count, the flag.}
\label{fig:review}
\end{figure}

\begin{figure}[htbp]\centering
\begin{tikzpicture}[font=\scriptsize]
  \node[draw=accent,fill=accent!10,rounded corners=1pt,align=center] (w) at (0,1.5) {write};
  \node[draw=red,fill=red!10,rounded corners=1pt,align=center] (r) at (3.4,1.5) {red-team};
  \node[draw=slate,fill=white,rounded corners=1pt,align=center] (f) at (6.8,1.5) {fix or\\ flag};
  \draw[->,ink,thick] (w)--(r) node[midway,above]{draft};
  \draw[->,ink,thick] (r)--(f) node[midway,above]{attacks};
  \draw[->,ink,dashed] (f.south) to[bend right=25] node[below]{repeat until flags are named} (w.south);
\end{tikzpicture}

\caption{The write-redteam-fix loop: repeat until every surviving flag has a name.}
\label{fig:loop}
\end{figure}

\section{From deck to memo}\label{sec:deck}

The transformation this book has been building toward: ``Sharpe 2.1'' becomes
``SR 2.1 in-sample / 0.4 out-of-sample, $N{=}40$ tried, flagged in quiet regime''---the
same research, the honest version.

\begin{figure}[htbp]\centering
\begin{tikzpicture}[font=\scriptsize]
  \node[draw=red,fill=red!10,rounded corners=1pt,align=left] at (0,1.6) {deck says:\\ ``Sharpe 2.1''};
  \node[draw=good,fill=goodbg,rounded corners=1pt,align=left] at (0,0.3) {memo says:\\ ``SR 2.1 IS / 0.4 OOS,\\ $N{=}40$ tried, quiet-regime flag''};
  \draw[->,red,thick] (3.0,1.7) to[bend left=15] node[above]{rewrite} (3.0,0.5);
\end{tikzpicture}

\caption{Deck to memo: the numbers are the same; the honesty is the rewrite.}
\label{fig:deck}
\end{figure}

\begin{example}\label{ex:deck}
\textbf{Why decks die.} A deck presents the best leaf of the forking tree; a memo presents
the tree.
\begin{solution}
Nobody can check ``Sharpe 2.1''---no data, no $N$, no OOS. The memo version can be
killed (weak OOS, tried 40 things) but can also be trusted: every field a reviewer
needs is present. The deliverable that survives is the one that volunteered its own
weaknesses.
\end{solution}
\end{example}

\begin{checkpoint}[Writeup audit]
\begin{itemize}
\item Name the six memo sections. (Claim, evidence, method, weaknesses, risks, repro.)
\item What are the three reproducibility pins? (Data hash, code commit, recorded seeds.)
\item What does a surviving flag mean? (A named weakness, not a failed check.)
\end{itemize}
\end{checkpoint}

\section{Recap}\label{sec:wurecap}

The memo is the last honest artifact in the chain: claim $\to$ evidence $\to$ trials
$\to$ weaknesses $\to$ risks $\to$ reproduction. Everything in this book---the harness,
the leaks, the validation, the metrics, the sizing, the risk---ends in a document
someone else can check. Write the one you'd want to review.

\begin{table}[htbp]\centering\small
\caption{Chapter toolbox.}
\begin{tabular}{@{}lll@{}}
\toprule
Concept & Rule & Gotcha\\
\midrule
Skeleton & 6 sections & order = reviewer's path\\
Claims & each carries evidence & evidence can refute\\
Repro & data+code+seeds & two of three is decoration\\
Red-team & run the attack list & the flag is the finding\\
Deck$\to$memo & honesty is the rewrite & same numbers, fuller truth\\
\bottomrule
\end{tabular}
\end{table}

\begin{example}[Linting the memo itself]\label{ex:memolint}
\textbf{Strategy.} Apply the book's own machinery to the writeup:
scan the findings section and require that every line containing a
number \emph{or} a verdict-word (``robust'', ``significant'',
``consistent'') carries a script reference in parentheses.

\textbf{Solution.} \texttt{memo\_lint.py} on a toy findings block:
\begin{verbatim}
memo lint: claims and their provenance
  line 1: OK   SMA(20,60) Sharpe +0.28 on 721 bars (run_sma.py, 5bp/2bp).
  line 2: OK   2024 slice Sharpe +4.08 (yearly_sharpe.py).
  line 3: MISS The edge is robust across regimes.
  line 4: OK   Kelly f* = 0.10 (kelly_sims.py).
3/4 claims cite a script
  FIX line 3: 'robust' is a claim --
     attach the regime table or delete the word
\end{verbatim}
The lint catches what a friendly reader misses: ``robust'' is a
claim, and it has no artifact behind it. The fix is mechanical ---
attach the regime table (ch12's per-year numbers) or delete the
word.

\textbf{Significance.} Reviewers trust claims they can re-run;
this lint is the automated form of the chapter's central rule.
Run it on your memo the way you run tests on code --- the findings
section is the only part of the book that \emph{is} its evidence.
\end{example}

\section{Common mistakes}
\begin{itemize}
\item \textbf{Claim without the number.} ``Works well'' is marketing, not research.
\item \textbf{Omitting the trial count.} The single most-omitted field in the industry.
\item \textbf{Results without seeds.} Unreproducible $\Rightarrow$ unverifiable
  $\Rightarrow$ anecdote.
\item \textbf{Hiding the OOS fold that lost.} The flag is the information.
\item \textbf{Repro instructions that are aspirational.} If a fresh clone can't run it,
  it doesn't reproduce.
\item \textbf{Concluding from the deck.} The deck is the pitch; the memo is the audit.
\item \textbf{Verdict words without artifacts.}
``Robust'', ``consistent'', ``significant'' are claims too ---
the lint flags them (Example~\ref{ex:memolint}). Attach the
table or delete the word.
\item \textbf{Evidence in the appendix only.} If the claim sits
in findings and its number sits in an appendix, the reader never
joins them --- keep claim and citation in the same paragraph.\end{itemize}

\section{Exercises}
\begin{exercise}[easy]
List the six memo sections.
\end{exercise}
\begin{solution}
Claim, evidence, method (incl.\ $N$), weaknesses, risks, reproduction.
\end{solution}

\begin{exercise}[easy]
Why does the claim come first?
\end{exercise}
\begin{solution}
It is the only part every reader reaches---and the part the evidence must answer.
\end{solution}

\begin{exercise}[easy]
Name the three reproducibility pins.
\end{exercise}
\begin{solution}
Data (hash), code (commit), seeds (recorded).
\end{solution}

\begin{exercise}[easy]
What is a ``flag'' in the red-team table?
\end{exercise}
\begin{solution}
A named weakness that survived the checklist---reported, not hidden.
\end{solution}

\begin{exercise}[easy]
Why is ``same seed $\to$ same numbers'' necessary but not sufficient?
\end{exercise}
\begin{solution}
It proves determinism, not correctness---the pins also need the data and code identities.
\end{solution}

\begin{exercise}
Rewrite ``Sharpe 2.1'' as a memo claim.
\end{exercise}
\begin{solution}
``SR 2.1 IS / 0.4 OOS on 720d BTC daily, $N{=}40$ variants tried, cost 5+2bp, flagged
in quiet regime.''
\end{solution}

\begin{exercise}
Draft the ``weaknesses'' section for the funding-carry memo.
\end{exercise}
\begin{solution}
Venue rule change, funding-sign flip in bear regimes, crowding of the carry, and the
500h sample covering mostly one regime.
\end{solution}

\begin{exercise}
Why does the evidence column get to say the claim is wrong?
\end{exercise}
\begin{solution}
The memo's job is truth, not advocacy---evidence that refutes is the highest-value kind.
\end{solution}

\begin{exercise}
Add a fourth pin to the three---what did you choose and why?
\end{exercise}
\begin{solution}
Environment (OS/library versions): same seed + different numpy can change the numbers.
\end{solution}

\begin{exercise}
Order the sections to match a hostile reviewer's read order.
\end{exercise}
\begin{solution}
Claim $\to$ evidence $\to$ method/trials $\to$ weaknesses $\to$ risks $\to$ repro---each
answers the next question a skeptic asks.
\end{solution}

\begin{exercise}
What separates the attack list from generic ``limitations'' prose?
\end{exercise}
\begin{solution}
Each attack is a checkable property (leak? costs? trials?) with a PASS/FLAG verdict,
not a paragraph of hedges.
\end{solution}

\begin{exercise}
A memo has 9/9 fields but no flag survived. Suspicious?
\end{exercise}
\begin{solution}
Yes---every real result has a boundary; a flagless memo usually means the red-team
wasn't run honestly.
\end{solution}

\begin{exercise}[hard]
Argue that the trial count field is the single most consequential in the memo, and
propose a convention for counting micro-choices.
\end{exercise}
\begin{solution}
Because $N$ converts a Sharpe into an inference (ch09); convention: declare every
variant evaluated, including those rejected before a formal test.
\end{solution}

\begin{exercise}[hard]
Design the ``repro'' section so a stranger can rerun the ch08 SMA result in under five
minutes.
\end{exercise}
\begin{solution}
`git clone; bash code/test\_08.sh'---the test greps the same numbers the chapter quotes;
hash of \texttt{data/btc\_daily.csv} is in the README.
\end{solution}

\begin{exercise}[hard]
Write the three-sentence version of this entire book.
\end{exercise}
\begin{solution}
Markets are measurable but not honest. Every tool that reports a number has a failure
mode that manufactures a number, so the research discipline is the chain: honest fills,
counted trials, purged folds, deflated bars, fractional sizing, and a memo that says
where it breaks.
\end{solution}

\begin{exercise}[easy]
Which two patterns does the memo lint search for before demanding
a script citation?
\end{exercise}
\begin{solution}
A numeric literal (decimal or signed percent) and a verdict word
(``robust'', ``significant'', ``consistent'', ``stable'',
``strong'') --- quantity or claim, either needs provenance.
\end{solution}
\begin{exercise}[medium]
Write the fix for ``The edge is robust across regimes'' without
deleting the claim.
\end{exercise}
\begin{solution}
``Per-year Sharpe ranged $-0.20$ to $+4.08$
(\texttt{yearly\_sharpe.py}) --- the edge lives in one regime''
--- keep the claim, attach the table, and let the numbers carry
the verdict.
\end{solution}
\begin{exercise}[hard]
Extend the lint to catch a subtler failure: a claim that cites a
script whose committed output no longer matches the number in the
memo. What metadata makes this checkable?
\end{exercise}
\begin{solution}
The script's pinned output hash / the data pin (sha256 of
\texttt{code/data/*}) plus the test that greps the quoted number
verbatim --- the memo cites the test file, not just the code, so
drift breaks the build.
\end{solution}

\section*{Notes}
This is the capstone the other fourteen chapters point at: none of the machinery matters
if the report doesn't survive a hostile read. The research that lasts is the research
that volunteered its weaknesses first---everything else gets red-teamed by someone else,
on their timeline, not yours.

\section*{Further reading}
\begin{itemize}
\item E.~Chan, ``What I Wish I'd Known About Quantitative Trading'' --- the practice:

\url{https://epchan.blogspot.com/}
\item L.~de Prado, ``The 7 Reasons Most Econometric Investments Fail'' --- the checklist:

\url{https://www.quantresearch.org/Lectures.htm}
\item ASA, ``Statistical inference'' (reproducibility standards):

\url{https://www.amstat.org/asa/files/pdfs/P-ValueStatement.pdf}
\item Project Jupyter (notebooks) --- the literate-research tool:

\url{https://jupyter.org/}
\item M.~Gardner, ``The Three Data Revolutions'' --- reproducibility culture:

\url{https://www.tandfonline.com/doi/full/10.1080/00031305.2023.2284320}
\end{itemize}

\begin{rulebox}[Watch instead]
\begin{itemize}
\item ``12. Risk Analysis'' --- MIT 18.S096 (80:32). What reviewers check; re-verified
2026-10-11.
\url{https://www.youtube.com/watch?v=qlytPllimpQ}
\item ``PyData pandas walkthrough'' --- McKinney (10:28). Reproducible workflow;
re-verified 2026-10-11.
\url{https://www.youtube.com/watch?v=_T8LGqJtuGc}
\end{itemize}
\end{rulebox}
